% GENERATED FILE. Edit canonical lessons, then run npm run generate:books.
% canonical-source-hash: fnv1a64:94edc4e7a8ee4100
% canonical-lessons: PA-C270-dhundhna, PA-C270-kharhna, PA-C270-langhauna, PA-C270-mil, PA-C270-maidan

\chapter{Searching for an Address}
\label{ch:pa-a2-searching-for-an-address}
\hlchaptermodality{270}

\begin{chapteropening}
\textbf{By the end of this chapter:} \emph{I can ask how far a place is and where to wait.}

Complete the last lesson of chapter 270: I can ask how far a place is and where to wait.
\end{chapteropening}

\section[\texorpdfstring{\pa{ਢੂੰਢਣਾ} (ḍhūṇḍhṇā)}{ਢੂੰਢਣਾ (ḍhūṇḍhṇā)}]{\pa{ਢੂੰਢਣਾ} (ḍhūṇḍhṇā) — to search for; to hunt for}

\label{lesson:PA-C270-dhundhna}



\begin{quote}
Before the new one: say the Punjabi for to draw near, then the Punjabi for a pothole.
\end{quote}

\subsection*{You'll want to know: \pa{ਢੂੰਢਣਾ}}
\textbf{\pa{ਢੂੰਢਣਾ}} — \emph{ḍhūṇḍhṇā} — \textquotedblleft{}to search for; to hunt for\textquotedblright{}.

\begin{quote}
\pa{ਮੈਂ} \pa{ਸ਼ਹਿਰ} \pa{ਵਿੱਚ} \pa{ਡਾਕਖ਼ਾਨਾ} \pa{ਢੂੰਢਦਾ} \pa{ਹਾਂ।}
\end{quote}

\begin{quote}
\emph{maiṁ shahir vich ḍākkhānā ḍhūṇḍhdā hāṁ.}
\end{quote}

\begin{quote}
I search the town for the post office.
\end{quote}

Past (I …): \pa{ਮੈਂ} \pa{ਥਾਣਾ} \pa{ਢੂੰਢਿਆ}, \emph{maiṁ thāṇā ḍhūṇḍhiā}.

\begin{cousinweb}
\textbf{\pa{ਢੂੰਢਣਾ}} (\emph{ḍhūṇḍhṇā}) is inherited from a reconstructed \emph{*ḍhūṇḍh}, \textquotedblleft{}seek\textquotedblright{}; CDIAL lists Punjabi \emph{ḍhū̃ḍhṇā}, \textquotedblleft{}to seek\textquotedblright{}, under it (CDIAL 6839).
\end{cousinweb}

\begin{grammarlens}[title={What you've built}]
One more word for this chapter.
\end{grammarlens}

\subsection*{Your turn}
\begin{itemize}
\raggedright
  \item \emph{Say it:} \emph{ḍhūṇḍhṇā}
  \item \emph{Say it:} \emph{ḍhūṇḍhṇā}, once more
  \item \emph{Say it:} \emph{maiṁ shahir vich ḍākkhānā ḍhūṇḍhdā hāṁ.} — I search the town for the post office.
  \item \emph{Recall:} say the Punjabi for to draw near, then the Punjabi for a pothole, then say \emph{ḍhūṇḍhṇā} again
\end{itemize}

\begin{tcolorbox}[breakable,colback=teal!4,colframe=teal!35!black,title={Before you move on}]
What does \pa{ਢੂੰਢਣਾ} mean? (\textquotedblleft{}To search for; to hunt for\textquotedblright{}.) Say it once more.
\end{tcolorbox}
\section[\texorpdfstring{\pa{ਖੜ੍ਹਨਾ} (khaṛhnā)}{ਖੜ੍ਹਨਾ (khaṛhnā)}]{\pa{ਖੜ੍ਹਨਾ} (khaṛhnā) — to stand; to stand waiting}

\label{lesson:PA-C270-kharhna}



\begin{quote}
Before the new one: say the Punjabi for a pothole, then the Punjabi for to search for.
\end{quote}

\subsection*{You'll want to know: \pa{ਖੜ੍ਹਨਾ}}
\textbf{\pa{ਖੜ੍ਹਨਾ}} — \emph{khaṛhnā} — \textquotedblleft{}to stand; to stand waiting\textquotedblright{}.

\begin{quote}
\pa{ਲੋਕ} \pa{ਕਤਾਰ} \pa{ਵਿੱਚ} \pa{ਖੜ੍ਹਦੇ} \pa{ਹਨ।}
\end{quote}

\begin{quote}
\emph{lok katār vich khaṛhde han.}
\end{quote}

\begin{quote}
People stand in the queue.
\end{quote}

Past (I …): \pa{ਮੈਂ} \pa{ਕਤਾਰ} \pa{ਵਿੱਚ} \pa{ਖੜ੍ਹਿਆ}, \emph{maiṁ katār vich khaṛhiā}.

\begin{cousinweb}
\textbf{\pa{ਖੜ੍ਹਨਾ}} (\emph{khaṛhnā}) goes back to Sanskrit \emph{khaḍaka}, \textquotedblleft{}erect; a post\textquotedblright{}, beside a reconstructed \emph{*khaḍati}, \textquotedblleft{}stands\textquotedblright{}; CDIAL lists Punjabi \emph{khaṛā}, \textquotedblleft{}standing erect\textquotedblright{}, under it (CDIAL 3784).
\end{cousinweb}

\begin{grammarlens}[title={What you've built}]
One more word for this chapter.
\end{grammarlens}

\subsection*{Your turn}
\begin{itemize}
\raggedright
  \item \emph{Say it:} \emph{khaṛhnā}
  \item \emph{Say it:} \emph{khaṛhnā}, once more
  \item \emph{Say it:} \emph{lok katār vich khaṛhde han.} — People stand in the queue.
  \item \emph{Recall:} say the Punjabi for a pothole, then the Punjabi for to search for, then say \emph{khaṛhnā} again
  \item \emph{Write it:} copy \textbf{\pa{ਖੜ੍ਹਨਾ}} once, checking the whole word
\end{itemize}

\begin{tcolorbox}[breakable,colback=teal!4,colframe=teal!35!black,title={Before you move on}]
What does \pa{ਖੜ੍ਹਨਾ} mean? (\textquotedblleft{}To stand; to stand waiting\textquotedblright{}.) Say it once more.
\end{tcolorbox}
\section[\texorpdfstring{\pa{ਲੰਘਾਉਣਾ} (laṅghāuṇā)}{ਲੰਘਾਉਣਾ (laṅghāuṇā)}]{\pa{ਲੰਘਾਉਣਾ} (laṅghāuṇā) — to see someone across; to help across}

\label{lesson:PA-C270-langhauna}



\begin{quote}
Before the new one: say the Punjabi for to search for, then the Punjabi for to stand.
\end{quote}

\subsection*{You'll want to know: \pa{ਲੰਘਾਉਣਾ}}
\textbf{\pa{ਲੰਘਾਉਣਾ}} — \emph{laṅghāuṇā} — \textquotedblleft{}to see someone across; to help across\textquotedblright{}.

\begin{quote}
\pa{ਮੈਂ} \pa{ਰੋਜ਼} \pa{ਬੱਚਿਆਂ} \pa{ਨੂੰ} \pa{ਸੜਕ} \pa{ਲੰਘਾਉਂਦਾ} \pa{ਹਾਂ।}
\end{quote}

\begin{quote}
\emph{maiṁ roz bachchiāṁ nū̃ saṛak laṅghāundā hāṁ.}
\end{quote}

\begin{quote}
Every day I help the children across the road.
\end{quote}

Past (I …): \pa{ਮੈਂ} \pa{ਬੱਚਿਆਂ} \pa{ਨੂੰ} \pa{ਪੁਲ} \pa{ਲੰਘਾਇਆ}, \emph{maiṁ bachchiāṁ nū̃ pul laṅghāiā}.

\begin{cousinweb}
\textbf{\pa{ਲੰਘਾਉਣਾ}} (\emph{laṅghāuṇā}) is the causative partner of \emph{laṅghṇā}, \textquotedblleft{}to pass\textquotedblright{}; CDIAL lists Punjabi \emph{laṅghṇā} under Sanskrit \emph{laṅghayati}, \textquotedblleft{}leaps over\textquotedblright{} (CDIAL 10905).
\end{cousinweb}

\begin{grammarlens}[title={What you've built}]
One more word for this chapter.
\end{grammarlens}

\subsection*{Your turn}
\begin{itemize}
\raggedright
  \item \emph{Say it:} \emph{laṅghāuṇā}
  \item \emph{Say it:} \emph{laṅghāuṇā}, once more
  \item \emph{Say it:} \emph{maiṁ roz bachchiāṁ nū̃ saṛak laṅghāundā hāṁ.} — Every day I help the children across the road.
  \item \emph{Recall:} say the Punjabi for to search for, then the Punjabi for to stand, then say \emph{laṅghāuṇā} again
  \item \emph{Write it:} copy \textbf{\pa{ਲੰਘਾਉਣਾ}} once, checking the whole word
\end{itemize}

\begin{tcolorbox}[breakable,colback=teal!4,colframe=teal!35!black,title={Before you move on}]
What does \pa{ਲੰਘਾਉਣਾ} mean? (\textquotedblleft{}To see someone across; to help across\textquotedblright{}.) Say it once more.
\end{tcolorbox}
\section[\texorpdfstring{\pa{ਮੀਲ} (mīl)}{ਮੀਲ (mīl)}]{\pa{ਮੀਲ} (mīl) — a mile}

\label{lesson:PA-C270-mil}



\begin{quote}
Before the new one: say the Punjabi for to stand, then the Punjabi for to see someone across.
\end{quote}

\subsection*{You'll want to know: \pa{ਮੀਲ}}
\textbf{\pa{ਮੀਲ}} — \emph{mīl} — \textquotedblleft{}a mile\textquotedblright{}.

\begin{quote}
\pa{ਸਕੂਲ} \pa{ਘਰ} \pa{ਤੋਂ} \pa{ਇੱਕ} \pa{ਮੀਲ} \pa{ਦੂਰ} \pa{ਹੈ।}
\end{quote}

\begin{quote}
\emph{sakūl ghar ton ikk mīl dūr hai.}
\end{quote}

\begin{quote}
The school is a mile from home.
\end{quote}

A masculine noun.

\begin{cousinweb}
\textbf{\pa{ਮੀਲ}} (\emph{mīl}) is borrowed from English \emph{mile}.
\end{cousinweb}

\begin{grammarlens}[title={What you've built}]
One more word for this chapter.
\end{grammarlens}

\subsection*{Your turn}
\begin{itemize}
\raggedright
  \item \emph{Say it:} \emph{mīl}
  \item \emph{Say it:} \emph{mīl}, once more
  \item \emph{Say it:} \emph{sakūl ghar ton ikk mīl dūr hai.} — The school is a mile from home.
  \item \emph{Recall:} say the Punjabi for to stand, then the Punjabi for to see someone across, then say \emph{mīl} again
  \item \emph{Write it:} copy \textbf{\pa{ਮੀਲ}} once, checking the whole word
\end{itemize}

\begin{tcolorbox}[breakable,colback=teal!4,colframe=teal!35!black,title={Before you move on}]
What does \pa{ਮੀਲ} mean? (\textquotedblleft{}A mile\textquotedblright{}.) Say it once more.
\end{tcolorbox}
\section[\texorpdfstring{\pa{ਮੈਦਾਨ} (maidān)}{ਮੈਦਾਨ (maidān)}]{\pa{ਮੈਦਾਨ} (maidān) — open ground; a playing field}

\label{lesson:PA-C270-maidan}



\begin{quote}
Before the new one: say the Punjabi for to see someone across, then the Punjabi for a mile.
\end{quote}

\subsection*{You'll want to know: \pa{ਮੈਦਾਨ}}
\textbf{\pa{ਮੈਦਾਨ}} — \emph{maidān} — \textquotedblleft{}open ground; a playing field\textquotedblright{}.

\begin{quote}
\pa{ਸਕੂਲ} \pa{ਦੇ} \pa{ਪਿੱਛੇ} \pa{ਇੱਕ} \pa{ਵੱਡਾ} \pa{ਮੈਦਾਨ} \pa{ਹੈ।}
\end{quote}

\begin{quote}
\emph{sakūl de pichchhe ikk vaḍḍā maidān hai.}
\end{quote}

\begin{quote}
There is a big playing field behind the school.
\end{quote}

A masculine noun.

\begin{cousinweb}
\textbf{\pa{ਮੈਦਾਨ}} (\emph{maidān}) is a loan, usually traced to Persian \emph{maidān}, \textquotedblleft{}open ground\textquotedblright{}.
\end{cousinweb}

\begin{grammarlens}[title={What you've built}]
One more word for this chapter.
\end{grammarlens}

\subsection*{Your turn}
\begin{itemize}
\raggedright
  \item \emph{Say it:} \emph{maidān}
  \item \emph{Say it:} \emph{maidān}, once more
  \item \emph{Say it:} \emph{sakūl de pichchhe ikk vaḍḍā maidān hai.} — There is a big playing field behind the school.
  \item \emph{Recall:} say the Punjabi for to see someone across, then the Punjabi for a mile, then say \emph{maidān} again
  \item \emph{Write it:} copy \textbf{\pa{ਮੈਦਾਨ}} once, checking the whole word
\end{itemize}

\begin{tcolorbox}[breakable,colback=teal!4,colframe=teal!35!black,title={Before you move on}]
What does \pa{ਮੈਦਾਨ} mean? (\textquotedblleft{}Open ground; a playing field\textquotedblright{}.) Say it once more.
\end{tcolorbox}
